\section{Failure Modes, Limitations, and Discussion}
\label{sec:discussion}

A method that improves every dataset by the same amount is suspicious; a method that fails in well-understood ways and improves the rest in a predictable pattern is trustworthy. PEM belongs to the second category. The theoretical account of Section~\ref{sec:theory} makes specific predictions about when entropy minimization on a converged checkpoint will help, when it will not help, and when it will actively hurt. This section documents the two principal failure modes we observe, discusses a third we anticipate but cannot evaluate from our data, addresses the limitations of our empirical setup, and positions the method in the broader landscape of post-hoc refinement.

\subsection{F1: Diffuse-entropy collapse}
\label{sec:disc:f1}

The first and most important failure mode occurs when the base checkpoint's residual entropy is not concentrated at the decision boundary. Formally, this is the regime $\rho_f(0.99) \lesssim 0.6$, in which a large fraction of the network's uncertainty lives on interior voxels whose features are not meaningfully bimodal. The supervised-only V-Net trained on the $12$ Pancreas-CT labeled volumes is the clearest example in our experiments: $\rho_f(0.99)\!=\!0.552$, and the mean $\ell_2$ distance from uncertain voxels to the predicted decision surface is $14.5$ voxels, compared to $5$--$11$ voxels on the SSL-converged checkpoints (Table~\ref{tab:entropy}). When PEM is applied to this checkpoint, the optimization halts within a single epoch: the collapse-protection safeguard that monitors training loss detects that the loss is increasing or flat rather than decreasing, and terminates the fine-tune to prevent damage.

The mechanism is the one predicted by Proposition~\ref{prop:gating} plus the cluster-assumption argument of Section~\ref{sec:theory:cluster}. At a diffuse-entropy voxel in the interior of a tissue, the feature distribution is not bimodal --- there is no partial-volume averaging, no scanner blur across a tissue interface, no anatomical ambiguity --- so the network's uncertainty at that voxel reflects its ignorance of local tissue texture rather than any genuine ambiguity in the input. The sign of the logit $z_i$ at such a voxel is essentially random: it reflects whichever class the network has latched onto without good reason. A PEM gradient step at this voxel pushes the logit further in whichever random direction it happens to be leaning, and does so with a step size proportional to $p_i(1{-}p_i)$ --- which is large precisely because the voxel is undecided. The result is \emph{commitment to an arbitrary choice}, which is the worst possible outcome: the network becomes more confident in a prediction it should not be making at all. Repeated over many voxels, this produces the loss increase that the collapse-protection safeguard detects.

\paragraph{Detecting F1 in advance.}
The $\rho_f(0.99)$ diagnostic is specifically designed to rule out F1 before any gradient step is taken. A single forward pass over the unlabeled training set suffices to compute $\rho_f$, and the threshold $\rho^\star\!=\!0.85$ separates the F1 regime from the PEM-suitable regime on every checkpoint we have tested. We therefore recommend that any deployment of PEM on a new checkpoint include a pre-flight measurement of $\rho_f(0.99)$ and a refusal to run if the measurement falls below $\rho^\star$. This is not an ad-hoc safeguard; it is the direct operational expression of the theoretical argument.

\paragraph{What F1 tells us about the applicability boundary.}
F1 is not a pathology of PEM as a loss; it is a property of the broader class of entropy-based refinement methods, and similar failure modes have been documented informally in the training-time entropy minimization literature for decades. What is new in our framing is that the regime boundary is captured by a single scalar computable from the unlabeled training set, and that the same scalar also predicts the magnitude of the expected improvement in the well-conditioned regime (Sec.~\ref{sec:theory:comparison}). The diagnostic doubles as a suitability check and as a rough effect-size estimator.

\subsection{F2: Vanishing-entropy diminishing returns}
\label{sec:disc:f2}

The second failure mode is less dramatic than the first and, importantly, is not a failure in the sense of producing a worse model. It is a diminishing-returns regime in which PEM continues to improve the base checkpoint but by an amount that shrinks toward zero as the base becomes more saturated. Formally, as the base checkpoint becomes more confident overall, $\rho_f(0.99)\!\to\!1$ from below \emph{and} the total residual entropy $\bar H$ shrinks. The soft entropy gradient of PEM scales with the mean Bernoulli variance $\bar V\!=\!\mathbb{E}_i[p_i(1{-}p_i)]$, which shrinks in step with $\bar H$; the magnitude of each PEM update therefore shrinks. The theoretical prediction of Section~\ref{sec:theory:comparison} is that, in this regime, a method whose gradients do not vanish with the total entropy budget --- specifically, pseudo-label fine-tuning, whose cross-entropy gradients are $O(1)$ regardless of confidence --- will become competitive with and eventually overtake PEM.

We observe exactly this on the LA 10\% configuration. The BCP base Dice is $89.40$, the highest of any configuration we evaluate; the corresponding $\bar H$ is correspondingly small. PEM improves Dice by $+0.64$ and PL-FT improves Dice by $+0.80$; their bootstrap $95\%$ confidence intervals overlap, and the two methods are statistically indistinguishable on this configuration by the per-case Wilcoxon test. On the two less-saturated configurations --- Pancreas-CT 20\% (baseline $82.89$) and LA 5\% (baseline $87.32$) --- PEM's gains exceed PL-FT's by $+0.62$ and $+0.34$ Dice respectively (see Table~\ref{tab:main}). The ordering across configurations is monotone in the base Dice in the direction the theory predicts: as the base model saturates, PEM's advantage over PL-FT shrinks and eventually reverses.

\paragraph{F2 is not a pathology; it is an exhaustion curve.}
Unlike F1, where PEM actively degrades the checkpoint, F2 simply means that there is less work for PEM to do. A fully saturated network has no residual entropy, hence no PEM gradient, hence no room to improve under this particular loss. The correct response to F2 is not to blame the method but to recognize the regime and consider alternative post-hoc refinement strategies --- most directly, pseudo-label fine-tuning with a conservative thresholding rule and connected-component filtering, whose $O(1)$ gradients continue to provide signal even when the soft entropy gradient has exhausted itself. A production system that automatically selects between PEM and PL-FT based on the measured $\rho_f$ and $\bar V$ would presumably achieve the maximum of the two curves across all saturation levels; we view this as a natural extension and leave it to future work.

\paragraph{An interpretation of the LA 10\% result.}
A reader accustomed to benchmark tables in which a single method dominates across all settings might read the LA 10\% row of Table~\ref{tab:main} as evidence that PEM is incomplete. We believe the opposite: the LA 10\% row is the empirical confirmation of the $p(1{-}p)$ prediction of Section~\ref{sec:theory:p1p}. The result is quantitatively consistent with the soft-gradient shrinkage argument, it is in the direction the theory predicts, and it does not invalidate the method's improvement on the other two settings. Reporting it honestly is one of the reasons we have confidence in the method's failure-mode characterization, and the observed tradeoff between soft and hard post-hoc gradients is a useful addition to the post-hoc refinement literature in its own right.

\subsection{F3: Unlabeled-set distribution shift}
\label{sec:disc:f3}

A third failure mode, which we anticipate but cannot evaluate from our current experiments, occurs when the unlabeled training set is drawn from a different distribution than the target evaluation distribution. If the unlabeled volumes come from a different scanner, a different patient population, a different contrast protocol, or a different anatomical variant, PEM will concentrate its refinement effort on whatever decision surface the base network has learned for \emph{that} distribution. If the target distribution differs, the refinement may push the decision surface in a direction that is locally optimal on the unlabeled set but globally suboptimal on the target set. This is not a new phenomenon in SSL --- it is the same pitfall that afflicts any method that uses unlabeled data as a proxy for the target distribution --- but it is worth flagging explicitly for PEM because the $\rho_f$ diagnostic does not detect it. A checkpoint can have high $\rho_f$ on the unlabeled set (well-localized residual uncertainty) and still be miscalibrated on a shifted target. In our experiments we use BCP's canonical splits, where labeled, unlabeled, and test volumes are drawn from the same distribution, so F3 does not arise. Evaluating PEM under deliberate unlabeled-set shift --- for example, PEM-fine-tuning on LA 5\% unlabeled volumes and evaluating on an independent LA test set from a different institution --- is a natural extension and would be the first serious test of PEM as a production refinement tool.

\subsection{Other failure modes we considered}
\label{sec:disc:other}

Two additional regimes deserve brief mention.

\paragraph{Extreme class imbalance.}
In domains where one class is rare (e.g.\ small lesions on a large background), the mean Bernoulli variance $\bar V$ is dominated by the majority class and may underestimate the amount of refinable signal on the minority-class boundary. In such cases $\rho_f(0.99)$ remains a reasonable indicator because it is normalized by total entropy and therefore reflects the distribution of uncertainty rather than its absolute magnitude, but the expected improvement from PEM on the minority-class Dice may be smaller than from PEM on the macro-average Dice. We do not encounter this regime on Pancreas-CT or LA, where the object classes are moderately sized relative to the background, and we have not evaluated PEM on small-lesion benchmarks such as BraTS or LiTS-tumor. The expected behavior is that $\rho_f(0.99)$ should still be a valid suitability check and that PEM should still improve the boundary between the lesion and the surrounding parenchyma, but the macro-average numbers may be more modest than on larger-organ tasks.

\paragraph{Base models that were never really converged.}
If the base SSL training is truncated early --- for instance, if the released checkpoint is from epoch $100$ of a recipe that normally runs for $300$ --- the network may not yet have entered the SSL convergence regime and may carry its residual entropy diffusely. We have not encountered this with the BCP checkpoints we use, all of which are fully trained, but a practitioner applying PEM to a half-trained checkpoint should expect the same diffuse-entropy collapse pattern as F1. The $\rho_f$ diagnostic detects this automatically.

\subsection{Limitations of our empirical evaluation}
\label{sec:disc:limits}

We discuss three limitations of the experimental evaluation that deserve acknowledgment.

\paragraph{Base-method coverage.}
Our empirical validation is restricted to BCP-family checkpoints because BCP is the only recent SSL method on Pancreas-CT and LA that releases public pretrained checkpoints. MOST~\cite{liu2024most} and DyCON~\cite{assefa2025dycon} provide training code but not checkpoints, and we could not faithfully reproduce their reported numbers within the time budget of this submission. Whether the $\rho_f$ concentration property and the resulting PEM improvement hold on MOST, DyCON, and other recent SSL checkpoints is an open question. Based on the theoretical argument of Section~\ref{sec:theory}, we expect the property to hold for any method that converges under a combination of supervised and consistency/pseudo-label signals, because the concentration is a property of the convergence rather than of the specific training loss; but this is a testable prediction we have not yet tested. Producing PEM-style results on MOST and DyCON, and characterizing the regime shift when moving between base methods, is the most important piece of follow-up work.

\paragraph{Dataset coverage.}
Our experiments cover two 3D medical segmentation tasks (Pancreas-CT organ segmentation and LA left-atrium segmentation) at three label fractions. Both tasks are large-organ binary segmentation with soft-tissue boundaries, which is the regime where the cluster-assumption-at-the-organ-surface argument of Section~\ref{sec:theory:cluster} is most directly applicable. We have not yet evaluated PEM on multi-class tasks such as ACDC~\cite{bernard2018acdc} or SynapseMulti, on small-lesion tasks such as BraTS~\cite{baid2021brats}, or on non-medical 3D segmentation tasks. The theoretical prediction is that the method should generalize to any setting where residual entropy concentrates at physically meaningful boundaries, but the quantitative magnitude of the improvement will depend on the concentration ratio $\rho_f$ and the total residual entropy budget, both of which are task-dependent. We view this as a limitation of scope rather than of the method and expect a follow-up paper to extend the empirical coverage substantially.

\paragraph{Learning-rate selection.}
PEM has one substantive hyperparameter, the learning rate, which we select once per dataset by $5$-fold cross-validation on the labeled training folds. The selected values ($5 \times 10^{-5}$ for Pancreas-CT, $1 \times 10^{-5}$ for LA) differ by a factor of five, reflecting the different degrees of saturation of the base checkpoints. An adaptive learning rate derived directly from the entropy histogram of the base checkpoint --- for instance, scaled inversely with $\bar V$ --- would remove this hyperparameter and would be more principled. We have experimented with such adaptive schemes but have not found one that matches the CV-selected rate on both datasets. A fully adaptive PEM is a natural extension and an attractive research target.

\subsection{Broader positioning and impact}
\label{sec:disc:broader}

PEM is a drop-in refinement step for any released SSL checkpoint on 3D medical image segmentation. It adds a small amount of compute to any training pipeline --- under five minutes per configuration on consumer hardware --- and is label-free, which means it can be applied by end users who do not have access to the original training data. In a research context, this makes it a useful tool for squeezing additional accuracy out of existing methods without retraining. In a deployment context, it means that a clinical site receiving a pretrained segmentation model from a research group can apply PEM to its own unlabeled institutional data as long as $\rho_f(0.99)$ exceeds the suitability threshold, with no exchange of patient labels. We have not benchmarked PEM in a clinical deployment setting, and the appropriate validation protocols for medical deployment (prospective studies, multi-site variation, failure-mode auditing beyond what we document here) are outside the scope of this paper.

We do not foresee negative societal impacts specific to PEM that do not already apply to any SSL segmentation method. The broader concerns around medical machine learning --- fairness across patient populations, robustness to distribution shift, interpretability, and the clinical appropriateness of any specific segmentation target --- apply to PEM exactly as they apply to the base SSL checkpoint it is refining. If anything, PEM's failure-mode characterization via $\rho_f$ offers a small amount of additional interpretability: a clinician using a PEM-refined checkpoint can, at least in principle, ask whether the refinement was well-posed on the institutional unlabeled data by measuring $\rho_f$ there, without needing labels.

\subsection{Directions for future work}
\label{sec:disc:future}

We close by listing the directions we find most promising.

\paragraph{Adaptive PEM from the entropy histogram.}
The most immediate extension is to remove the learning-rate hyperparameter by deriving it from $\bar V$ or from a histogram of $p(1{-}p)$ across voxels. A fully adaptive PEM would take any checkpoint and any unlabeled set as input and return a refined checkpoint with no tunable knobs.

\paragraph{A hybrid soft-and-hard post-hoc refinement.}
The F2 regime, where PL-FT becomes competitive with PEM, suggests that a combined loss --- soft entropy on low-confidence voxels, hard pseudo-label cross-entropy on high-confidence voxels beyond a carefully chosen threshold --- would dominate both individually. We have experimented with such combinations and found that they do, in fact, close the LA 10\% gap, but they introduce an additional hyperparameter (the threshold) and we have not found a principled way to set it from $\rho_f$ alone. A follow-up paper could make this hybrid rigorous.

\paragraph{PEM on non-BCP base methods.}
Producing PEM results on MOST, DyCON, MC-Net+, and any other publicly released 3D medical SSL checkpoint would immediately test the generality of the $\rho_f$ concentration property. If the property holds for any converged SSL checkpoint, PEM becomes a universal post-hoc refinement step for the field; if it does not, the characterization of which training recipes induce the concentration property and which do not is itself an interesting result.

\paragraph{PEM under unlabeled-set distribution shift.}
A controlled study of PEM's behavior when the unlabeled set and the target set come from different distributions would establish the method's robustness in production deployment. This is the most important open question for clinical use and is entirely within reach of a follow-up paper.

\paragraph{Beyond 3D segmentation.}
The cluster-assumption-at-the-physical-surface argument is specific to 3D medical segmentation, but it is not necessarily the only domain in which residual entropy concentrates at physically meaningful boundaries. 2D medical segmentation, natural-image semantic segmentation, and point-cloud segmentation are all candidates for the same style of regime analysis. The general question --- under what training recipes does a converged model leave its residual uncertainty in a spatially coherent, physically interpretable, refinable form? --- is a question about the geometry of convergence and is interesting well beyond the medical imaging application.

\subsection{Conclusion}
\label{sec:disc:conclusion}

We have introduced Post-hoc Entropy Minimization, a label-free, architecture-free, schedule-free fine-tuning step for converged SSL medical image segmentation checkpoints. The method is motivated by an empirical regularity: the residual prediction entropy of a converged SSL network is concentrated, to the tune of $89$--$95\%$, in a thin uncertain-voxel shell that coincides with the physical organ boundary. This concentration is what makes entropy minimization a refinement operator rather than a relocation operator in this regime, and is captured by a scalar label-free diagnostic $\rho_f(\tau)$ that also predicts when the method applies. The $p(1{-}p)$ factor of the binary-entropy gradient provides implicit boundary-shell gating, so no explicit mask or class-balance regularizer is needed; the same factor also predicts the diminishing-returns failure mode in the saturation limit. On the publicly released BCP checkpoints for Pancreas-CT 20\%, LA 5\%, and LA 10\%, PEM improves Dice by $+0.6$ to $+1.4$ and HD95 by $27$--$43\%$, in under $5.2$ minutes per configuration on a single GPU, without labels, pseudo-labels, or test-set hyperparameter tuning. Temperature scaling produces no Dice change on the same checkpoints, confirming that the improvement is not due to recalibration; the soft-gradient tradeoff against pseudo-label fine-tuning is quantitatively consistent with the theoretical prediction; and the two failure modes we characterize are exactly the ones the theory anticipates. We hope the $\rho_f$ diagnostic and the regime-shift framing are useful to the field independently of the specific method, and that the publicly released code and configurations will make it easy to apply PEM to any future SSL checkpoint on 3D medical image segmentation.
